\epigraph{\textit{``Indeed, brute force is a perfectly good technique in many cases; the real question is, can we use brute force in such a way that we avoid the worst-case behavior?''}}{--- \citeauthor{taocv3}, \citeyear{taocv3} \cite{taocv3}}

El enfoque de fuerza bruta consiste en generar de forma recursiva todas las posibles subsecuencias comunes entre las cadenas $s$ y $t$, seleccionando aquella de mayor longitud. A partir de esta subsecuencia común más larga (LCS), se reconstruyen las diferencias entre las cadenas comparando los bloques que quedan entre cada carácter coincidente.

El procedimiento recorre todas las combinaciones posibles avanzando en una de las dos cadenas o en ambas, y acumulando la subsecuencia cuando los caracteres coinciden. La diferencia entre las cadenas se obtiene después recorriendo $s$ y $t$ en paralelo, detectando fragmentos que no coinciden con la LCS.

Este método cumple con la definición clásica de fuerza bruta, ya que no evita ningún subproblema ni aplica poda o memoria intermedia.

\begin{algorithm}[H]
    \SetKwProg{myproc}{Procedure}{}{end}
    \SetKwFunction{BruteLCS}{BruteLCS}
    \SetKwFunction{GenerarDiferencias}{GenerarDiferencias}

    \footnotesize
    \DontPrintSemicolon

    \myproc{\BruteLCS{$s$, $t$, $i$, $j$, $actual$}}{
        \uIf{$i = |s|$ or $j = |t|$}{
            \If{$|actual| > |mejor\_subsecuencia|$}{
                $mejor\_subsecuencia \gets actual$\;
            }
            \Return\;
        }
        \uIf{$s[i] = t[j]$}{
            \BruteLCS{$s$, $t$, $i+1$, $j+1$, $actual + s[i]$}\;
        }
        \BruteLCS{$s$, $t$, $i+1$, $j$, $actual$}\;
        \BruteLCS{$s$, $t$, $i$, $j+1$, $actual$}\;
    }

    \myproc{\GenerarDiferencias{$s$, $t$}}{
        $mejor\_subsecuencia \gets ""$\;
        \BruteLCS{$s$, $t$, $0$, $0$, $""$}\;
        $lcs \gets mejor\_subsecuencia$\;
        Comparar $s$ y $t$ usando $lcs$ para extraer diferencias\;
    }

    \caption{Generación de LCS por fuerza bruta y reconstrucción de diferencias}
    \label{alg:brute_force_lcs}
\end{algorithm}

\textbf{Complejidad temporal:} En el peor de los casos, se generan todas las subsecuencias de $s$ y se comparan con todas las de $t$. Dado que una cadena de largo $n$ tiene $2^n$ subsecuencias, la complejidad de este algoritmo es $O(3^{\max(n,m)})$ al recorrer el espacio recursivo de decisiones.

\textbf{Complejidad espacial:} El algoritmo utiliza $O(n + m)$ espacio para representar la cadena actual acumulada y para la mejor subsecuencia encontrada. La pila de recursión tiene una profundidad máxima de $O(n + m)$.

Este enfoque no es viable para entradas grandes, pero cumple su propósito como referencia para análisis comparativo, y como demostración del paradigma de fuerza bruta.
